\documentclass[11pt,a4paper]{letter}
\usepackage[T1]{fontenc}
\usepackage[utf8]{inputenc}
\usepackage{lmodern}
\usepackage[margin=25mm]{geometry}
\usepackage[hidelinks]{hyperref}
\hypersetup{pdftitle={Cover Letter for Systems: Seeing the Process Through the Worker's Eyes},
  pdfauthor={Alessandro Berti and Wil M. P. van der Aalst}}

\signature{Alessandro Berti\\Corresponding author\\On behalf of both authors}
\address{Alessandro Berti\\Process and Data Science (PADS)\\RWTH Aachen University\\Aachen, Germany\\\href{mailto:a.berti@pads.rwth-aachen.de}{a.berti@pads.rwth-aachen.de}}
\date{\today}

\begin{document}
\begin{letter}{Guest Editors\\\textit{Systems}, MDPI\\Special Issue: Data-Driven Business Processes and Decision Optimization}
\opening{Dear Guest Editors,}

Please consider our research article, ``Seeing the Process Through the Worker's Eyes: Event Logs from Egocentric Factory Videos'', by Alessandro Berti and Wil M. P. van der Aalst, for the special issue ``Data-Driven Business Processes and Decision Optimization'' in \textit{Systems}.

Manual factory work often leaves few records in information systems. Our study presents a five-stage pipeline that uses foundation models to convert videos recorded from the worker's perspective into event logs. It creates a shared process and activity vocabulary for each factory and uses it to keep event labels consistent, without task-specific training or manual input annotations.

The resource contains 322 event logs with 10,698 events from 59 workers in six factories. The journal revision adds a formal event model, explicit measurement definitions, traceable examples, an analysis of missing time, and a portable workflow for calculations from frozen outputs. The existing author audit assesses agreement with generated transcripts; accuracy against video remains unmeasured. The logs quantify defect patching already described in a factory report and show how process-based categories can hide recorded transport. Timestamp inconsistencies, unequal observation windows, and shared textual evidence limit operational interpretation.

The paper fits the special issue's focus on process mining, process performance evaluation, and improvement in industrial systems. It connects human work with digital process analysis and provides evidence to guide process design decisions. The pipeline, prompts, derived data, and assessment code are publicly available.

The manuscript has not been published and is not under consideration for publication elsewhere. Both authors approve the manuscript and agree to its submission to \textit{Systems}. The manuscript includes declarations on author contributions, funding, ethics and consent, data availability, use of AI tools, and conflicts of interest.

Thank you for considering our submission.

\closing{Sincerely,}
\end{letter}
\end{document}
